\section{Introduction}
\label{sec:introduction}

In distributed communication systems, a central challenge lies in maintaining a secure but controlled flow of information between servers.
Although authorized entities require uninterrupted communication channels, adversarial agents must be denied access to any possible route that would allow information transfer.
This tension between enabling and restricting communication naturally motivates the study of graph structures that can simultaneously express both communication and isolation.
The interplay between these two objectives---connectivity and control---serves as the conceptual foundation of the present work.

Graph theory provides an elegant formal framework for representing such systems.
Vertices correspond to servers or agents, and edges represent direct communication links between them.
Within this model, ensuring communication between trusted servers while cutting off an opponent translates to identifying a set of edges that contains both a connecting path, which carries the communication, and a separating cut, which guarantees that every other route between the two servers uses at least one of the selected edges.
This hybrid structure, which captures both connection and disconnection properties at the same time, will be formally introduced as a \emph{cut-path}.
Such a structure allows the system to maintain communication within a trusted subnetwork while leaving an adversary no alternative route.

Several classical problems in network theory combine notions of connectivity and separation from different perspectives.
For example, the study of multi-terminal minimum cuts~\cite{GomoryHu1961}, of representations of all minimum cuts~\cite{NagamochiKameda1996}, and of small cuts and edge connectivity~\cite{Mehlhorn2017Certifying} explores the trade-off between maintaining essential communication routes and minimizing potential vulnerabilities.
Similarly, the interaction between \emph{shortest paths} and \emph{minimum cuts} has been investigated in the context of the \emph{max-flow min-cut duality}~\cite{Cormen2022}.
However, these formulations treat connectivity and separation as distinct or opposing properties, rather than unifying them within a single combinatorial object.
In contrast, the present work adopts a unified perspective through the concept of a \emph{cut-path}, which inherently embodies both connectivity and disconnection within the same set of edges.


The central problem studied in this paper, the \MinCutPath{} problem, seeks a smallest set of edges that simultaneously contains a path and a cut between two distinguished vertices.
This dual optimization objective reflects the balance between maintaining minimal communication infrastructure and enforcing maximal control over connectivity.
The problem was introduced in the author's master's thesis~\cite{Michel2025MinCutPath}, where several tractable cases were identified, but the complexity of the general problem was left open.
To the best of our knowledge, the problem has not been studied elsewhere.
We close this gap by settling the complexity of the problem and by presenting its basic structural properties.

Our main contributions are as follows.
First, we prove that the decision version of the \MinCutPath{} problem is NP-complete via a two-step reduction from \ThreeSAT{} through an intermediate problem called \SSP.
Second, we identify two graph classes in which the problem becomes tractable.
Namely, we show that for graphs of diameter two and for graphs with minimum cut-value at most two, the value of the minimum cut-path can be computed in polynomial time and is given by a simple closed formula.
Finally, we connect these results with random graph theory, demonstrating that dense Erdős–Rényi graphs almost surely fall into the tractable regime and that, in sparser random graphs, a simple algorithm yields an efficient average-case approximation.
The NP-completeness result is new; the results on special graph classes and on random graphs first appeared in~\cite{Michel2025MinCutPath}.

The remainder of this paper is structured as follows.
Section~\ref{sec:fundamentals} introduces the formal definitions, notation, and fundamental concepts used throughout the paper, including the definition of cut-paths and their basic properties.
The NP-completeness of the \MinCutPath{} problem is established in Section~\ref{sec:np-completeness} through a two-step reduction from \ThreeSAT{} via the intermediate \SSP{} problem.
Sections~\ref{sec:diameter-two} and~\ref{sec:cut-two} analyze two classes of graphs in which the problem becomes polynomially solvable: graphs of diameter two and graphs with minimum cut-value at most two.
Section~\ref{sec:random-graphs} investigates Erdős–Rényi random graphs, with emphasis on their diameter and connectivity, and derives an average-case approximation result for the \MinCutPath{} problem.
Finally, Section~\ref{sec:conclusion} concludes the paper and discusses directions for future research.

